\renewcommand{\theposttitle}{The Beginner's Mind and Finding Clarity}
\renewcommand{\thepostdate}{2022-07-08}
\vspace*{\fill}
\begin{center}
\includegraphics[width=\textwidth, height=0.55\textheight, keepaspectratio]{C:/Users/James/Documents/GitHub/VisualIdeasBook/data/images/88fed55bf7c04e9b80d33e30c706f162.png}
\end{center}
\vspace{2mm}
\begin{center}
\begin{minipage}{0.92\textwidth}
\small
The secret sauce to solving difficult problems is combining the efficiency of the expert with the openness of a beginners mind.There’s tremendous value in specialization and expertise.But there's also a hidden trap.To become an expert, you have to zoom in and learn deeply about a singular domain. The more you do something, the more efficient you get. But as you specialize and your efficiency improves, your attention gets narrower and you lose the ability to think laterally.Said another way, when you become an expert you sacrifice creativity for efficiency. The visual above celebrates the beginner's mind, or \textit{shoshin }(初心) as its known in Japanese Zen Buddhism. In the context of Zen, a beginner’s mind is like when you eat something for the very first time. Imagine that first bite. Its taste. Its texture. Its aroma...fully immersed in each moment of that experience. Having a beginner's mind is about: • snapping yourself out of auto-pilot mode• sustaining curiosity• and checking your assumptions about what's possible.

\textbf{DIVE DEEPER: }My favorite real world example of this is the the case of the scientist who discovered H. Pylori.
\end{minipage}
\end{center}
\vspace*{\fill}